\chapter{人工智能的历史与现状（代前言）}

人工智能的历史，大致可以分为三个阶段；每一个阶段都有其独特的技术路线、社会背景与兴衰轨迹。

\section{百花齐放的时代}

1956年的达特茅斯夏季研讨会标志着“人工智能”正式诞生。彼时，麦卡锡、明斯基、纽厄尔和西蒙等先驱带着极度乐观的情绪，开启了机器模仿人类智能的征程。

这个时代可谓“百花齐放”：符号推理、逻辑证明、神经网络（早期感知机）、进化算法、机器翻译等不同路径并行不悖。纽厄尔和西蒙的“逻辑理论家”程序能证明数学定理，塞缪尔的跳棋程序击败了人类选手，罗森布拉特的感知机则尝试直接模拟神经元学习。这些早期成功让学界和资助机构极度兴奋，美国国防部高级研究计划局（ARPA）等慷慨解囊，掀起第一次人工智能热潮。

兴起的原因十分清晰： 计算理论刚刚获得突破，研究者普遍相信，用几个简单定律就能复制智能；冷战背景下，自动翻译、实时情报分析等军事需求提供了充裕资金；而媒体和学术圈对“通用问题求解器”的过度宣传，催生了“十年内机器将能胜任所有人类工作”的期望泡沫。

然而，第一次热潮的衰落也同样迅速。首先，早期系统一旦超出微缩玩具世界，立刻遭遇“组合爆炸”——搜索空间随问题规模指数级增长，算法束手无策。机器翻译的荒唐错误使美国国家科学院在1966年发布ALPAC报告，认为其无望，导致经费大削。更具杀伤力的是，1969年明斯基与佩珀特出版《感知机》，严格证明了单层感知机连极为简单的异或问题都解决不了，并暗指多层网络同样前途渺茫\footnote{当然这个结论在现在看来是非常荒谬的。}，直接导致神经网络研究近乎绝迹。承诺不断落空，1973年英国莱特希尔报告严厉批评AI研究未能产生实效，全球资助急剧冰冻，第一次“AI之冬”降临。

这场衰落，本质上是当时算法能力、算力与理论工具与过度宏大的期望之间的巨大落差所致。

\section{符号主义的时代}

在寒冬中，一批研究者另辟蹊径，不再追求大一统的思考机器，转而聚焦于特定领域的专家知识，由此开启了符号主义的时代。所谓“符号主义”，是指以符号表示知识，并通过符号操作实现推理的人工智能方法。它的代表是“专家系统”，即将人类专家的经验和规则编码到计算机中，使其在特定领域内表现出类似专家的决策能力。

费根鲍姆等人倡导的知识工程，将人类专家的经验凝练成“如果-那么”规则，开发出能诊断血液感染病的MYCIN、推断分子结构的DENDRAL等专家系统。DEC公司的XCON系统每年能为企业节省数千万美元，成了商业活广告。资本迅速回流，各国政府争先追赶：日本抛出野心勃勃的“第五代计算机”计划，美欧纷纷跟进，以Lisp语言为核心的专用硬件和软件产业蓬勃一时，第二次人工智能热潮轰然展开。

这次热潮的根由，在于技术路线的务实转向。符号主义以物理符号系统假设为基，通过精心构造的知识库和推理机，在狭窄领域内展现了明确的经济价值，正好吻合了企业降本增效的迫切需求。

然而，泡沫再次迅速破灭。首要原因是知识获取瓶颈：把人类专家脑海中的隐性知识提取为显性规则，极其耗时耗力，且不同专家意见常相矛盾，规则库膨胀后变得异常脆弱、难以维护。其次，这些系统毫无常识，一旦越出狭窄边界就彻底失效，更不具备任何学习能力。压垮骆驼的最后一根稻草，是通用计算机性能飞速提升，专为Lisp设计的昂贵机器彻底丧失了市场优势；而日本五代机计划未能达成核心目标，宣告大规模符号主义路线的挫败。到1980年代末，产业溃缩，经费再度枯竭。

符号主义的衰落，根源于对知识静态、显性、可穷尽的幻想，未能面对真实世界知识模糊、开放且持续演化的本质。

\section{连接主义的时代}

连接主义的思想早已有之：模仿生物神经元的连接方式，设计“人工神经网络”，通过大量简单单元的并行计算实现复杂功能；但这套逻辑需要大量算力的支撑，那本《感知机》也给连接主义的相关研究投下了长长的阴影。

直到1986年，鲁梅尔哈特、辛顿等发表反向传播算法，才为多层神经网络的训练找到了可行路径。九十年代，神经网络因算力与数据稀缺、过拟合等问题，仍被支持向量机等统计学习模型盖过风头。

真正的引爆点出现在21世纪：2006年，辛顿提出深度信念网络的逐层预训练方法；2012年，亚历克斯·克里热夫斯基等利用GPU训练的深度卷积神经网络AlexNet，在ImageNet图像识别竞赛中以碾压性优势夺冠。从此，大数据、强算力、深度模型三者合力，在语音识别、计算机视觉、机器翻译、游戏博弈等一个又一个领域实现飞跃，“深度学习”一词家喻户晓，第三次人工智能热潮狂飙突进，至今潮头正高。

这次热潮的根基，是工程技术条件的质变。互联网的普及催生了海量标注数据，图形处理器和专用芯片的发展提供了空前的并行算力，算法细节的改进则让极深网络的训练成为可能。它用数据拟合函数，用算力换取智能，其效果之惊人，使产业界倾巢投入。

但从历史周期律来看，隐忧也在聚集。当前的连接主义模型大多是“数据饕餮”，依赖巨量算力，能耗惊人；以Transformer架构为首的模型决策如同黑箱，可解释性差，在医疗、司法等高风险领域处处受阻；系统大多脆弱，轻微扰动就可能导致识别错误；更根本的是，这种纯粹依赖关联统计的方法，缺乏因果推理与常识理解，智识层面仍显浅薄。若无法突破这些瓶颈，某一日喧嚣退去，人工智能也可能再次陷入寒冬。

\section{两种螺旋上升}

纵观人工智能的三段沉浮，一条螺旋上升的轨迹清晰可见。每一次热潮都不是凭空而来，而是在前一轮技术积累的废墟上，结合新的社会需求与技术条件，重新生长出来的。符号主义的兴起，脱胎于第一次热潮中对知识表示与推理的艰苦探索；连接主义的复兴，又深深受益于符号主义时代对计算模型和算法本质的透彻理解。今天，大语言模型的训练也远不只是“堆卡”那么简单——网络结构的精巧改进、基础算力设施（包括电力）的积累与飞跃，缺一不可。一项技术“做得出来”还是“做不出来”，说到底，是天时、地利、人和的共振。

但眼下有一种论调，说模型“训就完了”，似乎只要把数据灌进去、让显卡跑起来，智能便会自己涌现。我对这种看法一直心存疑虑。理论先于实践、体系严谨推导，这种做派确实是上世纪计算机科学的典型风格，在今天这个“先跑起来再说”的时代显得太慢，而现代模型又太大、不好解释，这种老古董的科研范式确实已不常见。然而，了解模型背后究竟在发生什么，这个条件下为什么能收敛、那个条件下又为什么怎么调都调不出来，仍然是真正理解人工智能不可绕开的关口。

初入人工智能的新手，往往容易死记硬背理论，公式背了一堆，却不知如何落地。等到有了一定的实践心得，又开始过度抬高工程的地位，觉得只要能把模型训出来就行，理论不过是纸上谈兵。殊不知这种撞大运式的做法，和买彩票并没有本质区别：个人的直觉与经验终有穷尽的一天，当模型怎么也训不好，乱调参数、乱改网络结构，最后的结果常常是“调了半天还是不行”，白白浪费大量时间与算力。这时候，很多人又会感慨书到用时方恨少，于是再次研究理论，看看前人有没有留下一些能够解释这类现象的线索。可理论学得多了又容易钻牛角尖，陷入“纸上谈兵”的窘境，于是又回到实践中去。

这种在理论与实践之间来来回回的过程，正是人工智能研究者最真实的日常。我们的知识边界与思维深度往往正是在这样一次次看似“回到原点”的折返中不断扩展。退一步说，就算暂时只看见螺旋，没感到上升，那也比原地踏步要强。

这本讲义取名《人工智能基础知识手册》，但我并不指望读者——尤其是初学者——第一遍就能全部消化。恰恰相反，真正领会这些理论为何必要、它们与代码之间怎样相互印证，恐怕非要亲身经历过几次螺旋不可：虽然我这里设计了一点PyTorch的实践内容，但实际上写出来的目前都是非常玩具的东西，和真正的模型差距很大：真要实践，还是得投身科研或参加工作，在项目中积累实践的经验。

至于这本讲义，处处看不懂，也完全没有关系——先看一遍，再看一遍，再看一遍。慢慢地，手感就来了，那些曾经干涩的公式，也会在某一天忽然间变得剔透起来。愿你在这螺旋中，走得远，也升得高。
